\documentclass[10pt]{article}
\usepackage[margin=0.98in]{geometry}
\usepackage{amsmath,amssymb,amsthm}
\usepackage[hidelinks]{hyperref}
\newtheorem{theorem}{Theorem}
\newcommand{\E}{\mathbb E}
\newcommand{\Pp}{\mathbb P}
\title{Full Counting Statistics of Neutral Emission in a Quantum Charged Throat}
\author{}
\date{}
\begin{document}
\maketitle
\begin{abstract}
In the deep-quantum, secular neutral-scalar evaporation model of a near-extremal charged black hole, the emission rate decreases as a power of the remaining energy and each emission removes a random fraction of that energy. The count therefore grows logarithmically in physical time rather than linearly. We prove the exact scaled cumulant generating function of the count and its large-deviation rate. The proof uses a change of the energy-loss distribution alone: it leaves the waiting-time law at a given energy unchanged. Uniform estimates on the residual energy show why this change of measure controls full count fluctuations and not merely a formal eigenvalue. The stationary count is sub-Poissonian and exceptionally quiet windows depend on whether the initial energy is fixed or drawn from the late ensemble. The results concern the single-channel continuum Markov model, not coherent multi-emission or the final evolution of an actual discharging black hole.
\end{abstract}

\section{Neutral throat input and physical time}
Let $E>0$ be the energy above extremality, and let $t$ be asymptotic proper time in the fixed external normalization. When the dominant open channel is the neutral scalar in the deep-quantum regime, the leading secular transition law has total rate
\begin{equation}
\lambda(E)=A E^{5/2},\qquad E\mapsto XE,\qquad
X\sim\mathrm{Beta}(3/2,2),\quad A>0.
\label{eq:physical}
\end{equation}
Inputs to this description include the near-throat Schwarzian dynamics, the appropriate greybody factors, and the secular evolution of energy populations. We take the rate and independent conditional fraction in \eqref{eq:physical} from the corresponding neutral continuum kernel; they are not new Hawking amplitudes. Biggs~\cite{Biggs} derives the energy master equation, its one-time evolution and late scaling law, in the context of quantum-corrected evaporation~\cite{IPU,BK}. Here the target is the joint path count, which is not determined by simply substituting the mean energy into a Poisson rate.

Fix a starting time $t_*>0$. Define the elapsed logarithmic time $s=\log(t/t_*)$ and the dimensionless remaining intensity $Z_s=AtE(t)^{5/2}$. The same physical events are counted on this clock: $N_S$ counts emissions between $t_*$ and $t_*e^S$. This is a deterministic change of the time label, not a different evaporation law or a level-dependent experimental timing protocol. Between emissions $dZ/ds=Z$; the intensity per unit $s$ is $Z$; and an event maps $Z$ to $ZV$ with $V=X^{5/2}$. For $b=3/5$,
\begin{equation}
m(u)=\E V^u=\frac{b}{(1+u)(b+u)},\qquad u>-b.
\label{eq:mellin}
\end{equation}
Equivalently $V$ is the product of independent uniform and $\mathrm{Beta}(b,1)$ variables. Thus the logarithmic energy decrements $Y=-\log V$ are a sum of independent exponentials of rates $1$ and $b$. The exact identity
\begin{equation}
\log Z_S-\log Z_0=S-\sum_{j=1}^{N_S}Y_j
\label{eq:identity}
\end{equation}
encodes the depletion mechanism: a count consumes an order-one amount of the log clock. The residual state term cannot be neglected without a bound, particularly under exponential tilting. We also use the known late scaling ensemble, expressed as $Z_0\overset{d}=WU$ with $W\sim\mathrm{Exp}(1)$ and $U\sim\mathrm{Beta}(b,1)$ independent~\cite{Biggs}; in the homogeneous logarithmic process it is stationary.

\section{Exact full-count growth}
\begin{theorem}
For a fixed $Z_0=z>0$, or for the stationary ensemble specified above, the real scaled cumulant generating function is
\begin{equation}
\lim_{S\to\infty}\frac{1}{S}\log\E e^{kN_S}
 =\alpha(k)=\frac{-(1+b)+\sqrt{(1-b)^2+4be^k}}{2},
 \qquad k\in\mathbb R.
\label{eq:scgf}
\end{equation}
The moment generating function is bounded above and below by finite strictly positive $k$-dependent multiples of $e^{\alpha(k)S}$ uniformly in $S$. In physical time the growth factor is $(t/t_*)^{\alpha(k)}$. The exponent is independent of the prefactor $A$ of the transition rate but depends on the Mellin law of the energy-loss marks.
\end{theorem}

\begin{proof}
The unique $\alpha>-b$ satisfying $e^k m(\alpha)=1$ is the stated root. Change the mark law to $\nu_{\alpha}(dv)=e^k v^\alpha\nu(dv)$, where $\nu$ is the original law of $V$. The jump rate at fixed $Z$ is unchanged, so the no-event likelihoods conditional on the path cancel. Nonexplosion on any finite interval follows from $Z_s\le z e^s$ at fixed start. The likelihood on a finite path and the resulting expectation identity are
\begin{align}
\frac{d\Pp_z^\alpha}{d\Pp_z}&=
e^{kN_S}\Bigl(\prod_{j=1}^{N_S}V_j\Bigr)^\alpha
=e^{kN_S-\alpha S}(Z_S/z)^\alpha,\\
\E_z e^{kN_S}&=e^{\alpha S}z^\alpha\E_z^\alpha Z_S^{-\alpha}.
\label{eq:change}
\end{align}
For the second line to control the exponent it remains to control both tails of $Z_S$ under the tilted law.

The generator on a smooth test $f$ is
\[
\mathcal L_\alpha f(z)=zf'(z)+z\int (f(zv)-f(z))\,\nu_\alpha(dv).
\]
Choose $p>\max(0,-\alpha)$ and $0<q<b+\alpha$, with $q>\alpha$ for $\alpha\ge0$. Such a $q$ exists for every finite $k$. For $U_p(z)= z^p+z^{-q}$ and $m_\alpha(u)=m(\alpha+u)/m(\alpha)$,
\[
\mathcal L_\alpha U_p=pz^p-(1-m_\alpha(p))z^{p+1}
-qz^{-q}+(m_\alpha(-q)-1)z^{1-q}.
\]
The strictly negative leading terms at infinity and zero show that $\mathcal L_\alpha U_p\le C-cU_p$ for $c>0$. Apply the integrating-factor Dynkin inequality on compact exit stops; then let the stops tend out to both ends and use nonexplosion, positivity and Fatou. This gives $\E_z^\alpha U_p(Z_S)\le e^{-cS}U_p(z)+C/c$. It bounds the terminal factor in \eqref{eq:change} above, and by uniform tightness inside a compact interval also bounds it below away from zero at each fixed start. The compact-uniform bounds also give a lower bound after averaging over a compact positive-probability set of stationary starts. The corresponding upper bound is integrable against the stationary law after multiplication by $z^\alpha$: all positive moments are finite, and the worst negative exponent $\alpha-q$ is greater than $-b$. This proves \eqref{eq:scgf} and the asserted finite factors.
\end{proof}

\section{Count fluctuations and quiet histories}
The count density $N_S/S$ obeys a large-deviation principle with rate
\begin{equation}
I(j)=
\begin{cases}
 j\log(jD/b)+(1+b-D)/2,&j>0,\\
 b,&j=0,\\
 +\infty,&j<0,
\end{cases}
\qquad
D=2j+\sqrt{(1-b)^2+4j^2}.
\end{equation}
Indeed $I$ is the Legendre transform of \eqref{eq:scgf}; differentiability on all real $k$ gives the interior bounds, positive exponential moments give exponential tightness, and the lower bound for relative neighborhoods of zero follows by approaching zero from positive densities. The law concentrates about $j=b/(1+b)=3/8$, so the leading number of events is $3\log(t/t_*)/8$ in the indefinitely continued ideal process.

An independent variance check uses the physical depletion identity rather than assuming that derivatives commute with the cumulant limit. In the stationary ensemble the count intensity has finite expectation, and $\log Z$ has finite second moment. With $\mu_Y=1+1/b$, $\sigma_Y^2=1+1/b^2$, the compensated mark martingale $M_S=\sum_{j\le N_S}(Y_j-\mu_Y)$ has mean zero and variance $\sigma_Y^2\E N_S$, since each new mark is independent of its past conditional history. Equation \eqref{eq:identity} yields $\E N_S=S/\mu_Y$ and $\mu_Y N_S=S-M_S-(\log Z_S-\log Z_0)$. The endpoint term has bounded variance; its covariance with $M_S$ is $O(\sqrt S)$. Hence the asymptotic variance-to-mean ratio is $\sigma_Y^2/\mu_Y^2=(1+b^2)/(1+b)^2=17/32$. The sub-Poissonian result follows here from a classical stochastic depletion process with quantum-corrected input; it is not evidence of coherence by itself.

Exactly zero events require additional care. For a fixed start the intensity keeps rising as $ze^s$ along the no-event path, giving $\Pp_z(N_S=0)=e^{-z(e^S-1)}$. At stationary start,
\[
\Pp(N_S=0)=b\int_0^1\frac{v^{b-1}}{1+(e^S-1)v}\,dv
\sim\frac{\pi b}{\sin(\pi b)}e^{-bS}.
\]
Late time averaging allows rare very low initial intensities, explaining the different quiet-window behavior. Neither statement contradicts the large-deviation rate of neighborhoods of count density zero.

\section{Interpretation and range of validity}
The mechanism is not specific to the scalar's square-root formula. In the general single-channel kernel of Biggs~\cite[Eq. (34)]{Biggs}, the transition density from $E$ to $E'\in(0,E)$ is proportional to $(E')^{1/2}(E-E')^a$, where a fixed integer $a\ge0$ labels the power of the dominant neutral channel. Writing $\gamma=a+3/2$, the rate is $A_a E^\gamma$ and $X=E'/E$ has law $\mathrm{Beta}(3/2,a+1)$. On the same physical log clock, $V=X^\gamma$ and
\[
\E V^u=\prod_{j=0}^a\frac{b_j}{b_j+u},\qquad b_j=\frac{3/2+j}{\gamma}.
\]
Thus the logarithmic decrement is a sum of exponentials with rates $b_0,\ldots,b_a$. At fixed start, the same endpoint proof goes through with $b$ replaced by $b_0$ in the restriction on negative moments. The exact count exponent is the unique real root $\alpha>-b_0$ of $\prod_{j=0}^a(1+\alpha/b_j)=e^k$. This yields one dominant-channel law for each prescribed kernel (the scalar is $a=1$, photon $a=3$, and the di-photon transition $a=7$). In the last case the count is events, not photons, and the formula is not for concurrent channels or crossovers between them.

The $S\to\infty$ result is a statement about the homogeneous effective process when continued throughout arbitrarily large $t/t_*$. A real measurement can test only windows for which the remnant stays in the deep-quantum near-extremal continuum and the chosen neutral scalar channel dominates. In particular spectral discreteness, coherent multi-quantum emission, other species or angular momentum, and charge discharge can end that window. No exchange of the large log-window limit with any of those physical limits is asserted. The result does not prescribe a physical endpoint, identify a microscopic unitary evaporation model, or rely on changing the physical clock as the energy falls.

\begin{thebibliography}{9}
\bibitem{Biggs}
A.~Biggs, \emph{Following the state of an evaporating charged black hole into the quantum gravity regime}, \href{https://arxiv.org/abs/2503.02051}{arXiv:2503.02051}, especially Sec.~3.2.
\bibitem{IPU}
L.~V.~Iliesiu, G.~Penington and M.~Usatyuk, \emph{The evaporation of charged black holes}, JHEP \textbf{01} (2026) 109, \href{https://doi.org/10.1007/JHEP01(2026)109}{doi:10.1007/JHEP01(2026)109}.
\bibitem{BK}
Y.~Bai and M.~Korwar, \emph{Near-extremal charged black holes: greybody factors and evolution}, JHEP \textbf{03} (2023) 151, \href{https://doi.org/10.1007/JHEP03(2023)151}{doi:10.1007/JHEP03(2023)151}.
\end{thebibliography}
\end{document}
